\chapter{Glossary of terms}
\label{ch:glossary2}

\begin{aside}
A reference for the vocabulary used across this book.
Alphabetical, brief.
\end{aside}

\textbf{Alt-screen.} The terminal's separate buffer for
full-screen apps. \maya{} enters on startup.

\textbf{ANSI escape.} A control sequence that changes
terminal state. Always starts with \code{ESC} (\code{0x1B}).

\textbf{Awaiter.} An object that suspends a coroutine and
resumes it later with a value.

\textbf{Back buffer.} The off-screen canvas where painting
happens before diff and emit.

\textbf{Bracketed paste.} Mode where the terminal wraps
pasted bytes in markers so the app can distinguish from
typed input.

\textbf{Canvas.} The 2-D grid of cells representing the
screen.

\textbf{Cell.} One column of one row. \maya{} packs each
into 64 bits.

\textbf{Cluster.} A grapheme cluster: the user-perceived
``character''.

\textbf{Coroutine.} A function that can pause and resume.
\cpp{}20 added the syntax.

\textbf{CSI-u.} Modern keyboard encoding that gives every
key a unique escape.

\textbf{DEC 2026.} Synchronised output mode; brackets a
frame so the terminal applies it atomically.

\textbf{Diff.} Compare two canvases; produce a minimal
list of changes.

\textbf{Effect.} A signal subscriber that produces side
effects when its dependencies change.

\textbf{Element.} A node in the UI tree. Layout,
positioning, paint.

\textbf{Front buffer.} The canvas representing what's on
screen.

\textbf{Grapheme.} See cluster.

\textbf{Memo.} A computed signal whose value is cached
until inputs change.

\textbf{Microtask.} A deferred function scheduled to run
soon, not immediately.

\textbf{OSC-8.} Hyperlink escape sequence.

\textbf{OSC-52.} Clipboard escape sequence.

\textbf{Owner.} Lifetime parent for signals and effects.

\textbf{Painter.} The API for writing into the back buffer.

\textbf{Pty.} Pseudo-terminal: bidirectional fd pair that
mimics a real terminal.

\textbf{Reactive.} State + automatic recomputation of
derived values.

\textbf{Reconcile.} Compare new and old element trees,
compute minimal updates.

\textbf{Render.} A frame: layout, paint, diff, emit.

\textbf{SGR.} Set Graphics Rendition: ANSI for colour and
attributes.

\textbf{Signal.} Reactive container for a value.

\textbf{Sixel.} DEC's pixel encoding for terminal graphics.

\textbf{Sprite.} A cached cell array for fast blit.

\textbf{TUI.} Text User Interface. The category.

\textbf{Truecolour.} 24-bit RGB; supported by modern
terminals.

\textbf{Variation selector.} Code point that picks emoji
or text presentation.

\textbf{Virtualised list.} List that renders only visible
rows.

\textbf{Widget.} A reusable element type.

\textbf{XDG.} Cross-Desktop Group: standard directory
conventions.

\textbf{ZWJ.} Zero-width joiner; emoji sequence connector.
